% Options for packages loaded elsewhere
% Options for packages loaded elsewhere
\PassOptionsToPackage{unicode}{hyperref}
\PassOptionsToPackage{hyphens}{url}
\PassOptionsToPackage{dvipsnames,svgnames,x11names}{xcolor}
%
\documentclass[
  french,
  11pt,
  a4paper,
  twoside,
  article,
  titlepage=false]{memoir}
\usepackage{xcolor}
\usepackage{amsmath,amssymb}
\setcounter{secnumdepth}{-\maxdimen} % remove section numbering
\usepackage{iftex}
\ifPDFTeX
  \usepackage[T1]{fontenc}
  \usepackage[utf8]{inputenc}
  \usepackage{textcomp} % provide euro and other symbols
\else % if luatex or xetex
  \usepackage{unicode-math} % this also loads fontspec
  \defaultfontfeatures{Scale=MatchLowercase}
  \defaultfontfeatures[\rmfamily]{Ligatures=TeX,Scale=1}
\fi
\usepackage{lmodern}
\ifPDFTeX\else
  % xetex/luatex font selection
\fi
% Use upquote if available, for straight quotes in verbatim environments
\IfFileExists{upquote.sty}{\usepackage{upquote}}{}
\IfFileExists{microtype.sty}{% use microtype if available
  \usepackage[]{microtype}
  \UseMicrotypeSet[protrusion]{basicmath} % disable protrusion for tt fonts
}{}
\makeatletter
\@ifundefined{KOMAClassName}{% if non-KOMA class
  \IfFileExists{parskip.sty}{%
    \usepackage{parskip}
  }{% else
    \setlength{\parindent}{0pt}
    \setlength{\parskip}{6pt plus 2pt minus 1pt}}
}{% if KOMA class
  \KOMAoptions{parskip=half}}
\makeatother
% Make \paragraph and \subparagraph free-standing
\makeatletter
\ifx\paragraph\undefined\else
  \let\oldparagraph\paragraph
  \renewcommand{\paragraph}{
    \@ifstar
      \xxxParagraphStar
      \xxxParagraphNoStar
  }
  \newcommand{\xxxParagraphStar}[1]{\oldparagraph*{#1}\mbox{}}
  \newcommand{\xxxParagraphNoStar}[1]{\oldparagraph{#1}\mbox{}}
\fi
\ifx\subparagraph\undefined\else
  \let\oldsubparagraph\subparagraph
  \renewcommand{\subparagraph}{
    \@ifstar
      \xxxSubParagraphStar
      \xxxSubParagraphNoStar
  }
  \newcommand{\xxxSubParagraphStar}[1]{\oldsubparagraph*{#1}\mbox{}}
  \newcommand{\xxxSubParagraphNoStar}[1]{\oldsubparagraph{#1}\mbox{}}
\fi
\makeatother


\usepackage{longtable,booktabs,array}
\newcounter{none} % for unnumbered tables
\usepackage{calc} % for calculating minipage widths
% Correct order of tables after \paragraph or \subparagraph
\usepackage{etoolbox}
\makeatletter
\patchcmd\longtable{\par}{\if@noskipsec\mbox{}\fi\par}{}{}
\makeatother
% Allow footnotes in longtable head/foot
\IfFileExists{footnotehyper.sty}{\usepackage{footnotehyper}}{\usepackage{footnote}}
\makesavenoteenv{longtable}
\usepackage{graphicx}
\makeatletter
\newsavebox\pandoc@box
\newcommand*\pandocbounded[1]{% scales image to fit in text height/width
  \sbox\pandoc@box{#1}%
  \Gscale@div\@tempa{\textheight}{\dimexpr\ht\pandoc@box+\dp\pandoc@box\relax}%
  \Gscale@div\@tempb{\linewidth}{\wd\pandoc@box}%
  \ifdim\@tempb\p@<\@tempa\p@\let\@tempa\@tempb\fi% select the smaller of both
  \ifdim\@tempa\p@<\p@\scalebox{\@tempa}{\usebox\pandoc@box}%
  \else\usebox{\pandoc@box}%
  \fi%
}
% Set default figure placement to htbp
\def\fps@figure{htbp}
\makeatother



\ifLuaTeX
\usepackage[bidi=basic,shorthands=off]{babel}
\else
\usepackage[bidi=default,shorthands=off]{babel}
\fi
\ifLuaTeX
  \usepackage{selnolig} % disable illegal ligatures
\fi


\setlength{\emergencystretch}{3em} % prevent overfull lines

\providecommand{\tightlist}{%
  \setlength{\itemsep}{0pt}\setlength{\parskip}{0pt}}



 


\makepagestyle{style-td1}
\makeevenhead{style-td1}{\sffamily\small Intégration et Probabilités}{\sffamily B}{\sffamily\small CC III MIASHS 1 \& 2}
\makeoddhead{style-td1}{\sffamily\small Intégration et Probabilités}{\sffamily B}{\sffamily\small CC III MIASHS 1 \& 2}
\makeevenfoot{style-td1}{\sffamily\small MA1AY010}{\thepage}{\sffamily\small L3 MIASHS}
\makeoddfoot{style-td1}{\sffamily\small L3 MIASHS}{\thepage}{\sffamily\small MA1AY010}
\pagestyle{style-td1}
\makeatletter
\@ifpackageloaded{tcolorbox}{}{\usepackage[skins,breakable]{tcolorbox}}
\@ifpackageloaded{fontawesome5}{}{\usepackage{fontawesome5}}
\definecolor{quarto-callout-color}{HTML}{909090}
\definecolor{quarto-callout-note-color}{HTML}{0758E5}
\definecolor{quarto-callout-important-color}{HTML}{CC1914}
\definecolor{quarto-callout-warning-color}{HTML}{EB9113}
\definecolor{quarto-callout-tip-color}{HTML}{00A047}
\definecolor{quarto-callout-caution-color}{HTML}{FC5300}
\definecolor{quarto-callout-color-frame}{HTML}{acacac}
\definecolor{quarto-callout-note-color-frame}{HTML}{4582ec}
\definecolor{quarto-callout-important-color-frame}{HTML}{d9534f}
\definecolor{quarto-callout-warning-color-frame}{HTML}{f0ad4e}
\definecolor{quarto-callout-tip-color-frame}{HTML}{02b875}
\definecolor{quarto-callout-caution-color-frame}{HTML}{fd7e14}
\makeatother
\makeatletter
\@ifpackageloaded{caption}{}{\usepackage{caption}}
\AtBeginDocument{%
\ifdefined\contentsname
  \renewcommand*\contentsname{Table des matières}
\else
  \newcommand\contentsname{Table des matières}
\fi
\ifdefined\listfigurename
  \renewcommand*\listfigurename{Liste des Figures}
\else
  \newcommand\listfigurename{Liste des Figures}
\fi
\ifdefined\listtablename
  \renewcommand*\listtablename{Liste des Tables}
\else
  \newcommand\listtablename{Liste des Tables}
\fi
\ifdefined\figurename
  \renewcommand*\figurename{Figure}
\else
  \newcommand\figurename{Figure}
\fi
\ifdefined\tablename
  \renewcommand*\tablename{Table}
\else
  \newcommand\tablename{Table}
\fi
}
\@ifpackageloaded{float}{}{\usepackage{float}}
\floatstyle{ruled}
\@ifundefined{c@chapter}{\newfloat{codelisting}{h}{lop}}{\newfloat{codelisting}{h}{lop}[chapter]}
\floatname{codelisting}{Listing}
\newcommand*\listoflistings{\listof{codelisting}{Liste des Listings}}
\usepackage{amsthm}
\theoremstyle{definition}
\newtheorem{exercise}{Exercice}[section]
\theoremstyle{remark}
\AtBeginDocument{\renewcommand*{\proofname}{Preuve}}
\newtheorem*{remark}{Remarque}
\newtheorem*{solution}{Solution}
\newtheorem{refremark}{Remarque}[section]
\newtheorem{refsolution}{Solution}[section]
\makeatother
\makeatletter
\makeatother
\makeatletter
\@ifpackageloaded{caption}{}{\usepackage{caption}}
\@ifpackageloaded{subcaption}{}{\usepackage{subcaption}}
\makeatother
\usepackage{bookmark}
\IfFileExists{xurl.sty}{\usepackage{xurl}}{} % add URL line breaks if available
\urlstyle{same}
% fallback for those not using the hyperref driver hyperxmp:
\makeatletter
\@ifundefined{xmpquote}{\newcommand{\xmpquote}[1]{#1}}{}
\makeatother
\hypersetup{
  pdflang={fr},
  colorlinks=true,
  linkcolor={blue},
  filecolor={Maroon},
  citecolor={Blue},
  urlcolor={Blue},
  pdfcreator={LaTeX via pandoc}}


\author{}
\date{2026-06-03}
\begin{document}
\frontmatter

\counterwithout{exercise}{section}


\mainmatter
\begin{tcolorbox}[enhanced jigsaw, arc=.35mm, bottomrule=.15mm, breakable, colback=white, colframe=quarto-callout-note-color-frame, left=2mm, leftrule=.75mm, opacityback=0, rightrule=.15mm, toprule=.15mm]

\vspace{-3mm}\textbf{CC III/IV :}\vspace{3mm}

\begin{itemize}
\tightlist
\item
  Licence MIASHS
\item
  15 Décembre 2025 Durée : 1 heure 15
\item
  \textbf{Intégration et Probabilités}
\end{itemize}

\end{tcolorbox}

\begin{exercise}[]\protect\hypertarget{exr-1-a}{}\label{exr-1-a}

~

\end{exercise}

\begin{tcolorbox}[enhanced jigsaw, arc=.35mm, bottomrule=.15mm, breakable, colback=white, colframe=quarto-callout-note-color-frame, left=2mm, leftrule=.75mm, opacityback=0, rightrule=.15mm, toprule=.15mm]

Dans cet exercice \(\mathcal{N}(\mu, \Sigma)\) désigne la loi d'un
vecteur gaussien d'espérance \(\mu \in \mathbb{R}^d\) et de matrice de
covariance \(\Sigma\) à \(d\) lignes et \(d\) colonnes.

\end{tcolorbox}

Dans cet exercice,
\(X = (X_1, X_2, X_3)^{\top} \sim \mathcal{N}(0, \Sigma)\) avec

\[
\Sigma = 
  \begin{pmatrix} 
  1 & 1 & 1 \\ 
  1 & 2 & 2 \\ 
  1 & 2 & 3 \\ 
  \end{pmatrix}
\]

\begin{enumerate}
\def\labelenumi{\arabic{enumi}.}
\tightlist
\item
  La loi de \(X\) admet-elle une densité par rapport à la mesure de
  Lebesque sur \(\mathbb{R}^3\) ? Si oui, donner cette densité.
\item
  Calculer la transformée de la Laplace de la loi \(X_3 - X_1\).
\item
  Calculer la densité conditionnelle de \(X_3\) sachant \(X_1=x_1\) et
  \(X_2=x_2\) ?
\item
  Calculer l'espérance et la variance de \(X_1^2 + X_2^2\)
\end{enumerate}

\begin{exercise}[]\protect\hypertarget{exr-3-a}{}\label{exr-3-a}

~

\end{exercise}

Dans cet exercice, \(\lambda\) prend ses valeurs sur \(]0, \infty)\),
\(Z_{\lambda}\) est une variable aléatoire à valeur dans \(\mathbb{N}\).
Sa loi est définie par sa fonction génératrice

\[G_\lambda(s) = \exp\left(\lambda \frac{s -1 }{1 - (1-p)s} \right)\]

pour \(s \in [0,1]\), \(\lambda>0\) et \(p\in (0,1)\)

\begin{enumerate}
\def\labelenumi{\arabic{enumi}.}
\item
  Calculer la probabilité que \(Z\) prenne les valeurs, \(0\) et \(1\)
\item
  Calculer l'espérance de \(Z\) (si finie)
\item
  Calculer la variance de \(Z\) (si finie)

  \emph{Maintenant \(\lambda=1\), et
  \(X_1, \ldots, X_n, \ldots \sim_{\text{i.i.d.}} G_1\) (distribuées
  selon la loi définie par \(G_1\)), et \(S_n = \sum_{i=1}^n X_i\).}
\item
  Calculer la fonction génératrice de la loi de \(S_n\)
\item
  Étudier la convergence en probabilité de la suite \((S_n/n)_n\)
\end{enumerate}

\newpage{}

\begin{exercise}[]\protect\hypertarget{exr-5-a}{}\label{exr-5-a}

~

\end{exercise}

\begin{tcolorbox}[enhanced jigsaw, arc=.35mm, bottomrule=.15mm, breakable, colback=white, colframe=quarto-callout-note-color-frame, left=2mm, leftrule=.75mm, opacityback=0, rightrule=.15mm, toprule=.15mm]

Rappel: \emph{La loi exponentielle de paramètre d'intensité
\(\lambda>0\) a pour densité
\(x \mapsto \mathbb{I}_{x>0} \lambda \exp(-\lambda x)\).}

\end{tcolorbox}

Dans cet exercice \(Z_1, Z_2, Z_{1,2}\) sont trois variables
indépendantes distribuées selon des loi exponentielles d'intensités
respectives \(1, 1, 2\).

On définit, \(X = \min(Z_1, Z_{1,2})\) et \(Y =  \min(Z_2, Z_{1,2})\).

\protect\phantomsection\label{marshall-olkin-ques}
\begin{enumerate}
\def\labelenumi{\arabic{enumi}.}
\tightlist
\item
  Caractérisez la loi de \(X\) (fonction de répartition, densité,
  fonction de survie).
\item
  Caractérisez la loi de \(\min(X, Y)\)
\item
  Quelle est la probabilité que \(X=Y\) ?
\item
  \(X\) et \(Y\) sont-elles indépendantes ?
\item
  La loi jointe de \((X, Y)\) admet-elle une densité par rapport à la
  mesure de Lebesgue sur \(\mathbb{R}^2\) ? Si oui, calculer la densité.
\item
  Calculer \(P\{ X > a \text{ et } Y > b\}\) pour
  \(a, b ∈ ]0,\infty)^2\)
\item
  Calculer
  \(P\{ X > a + s \text{ et } Y > b +s \mid X >a \text{ et } Y>b\}\)
  pour \(a, b, s ∈ ]0,\infty)^2\)
\end{enumerate}


\backmatter


\end{document}
